\section{Introduction}
\label{sec:introduction}

The deployment of autonomous multi-agent systems into high-stakes operational environments has exposed a structural failure mode that existing safety architectures are not equipped to address: the problem of \textit{autonomous drift}. In brief, drift occurs when an agent continues to execute, spawn descendants, and take real-world actions with no traceable authorization path to a human principal. Unlike a system crash or a clear runtime error, drift is not self-evidently broken. The agent persists, optimizes, propagates — but the chain of delegation that warranted its agency has been severed, or was never architecturally established. The result is an agent exercising genuine consequential authority that no human has recently authorized or can reliably revoke.

The stakes of drift are not hypothetical. In recursive spawning environments — where agents are authorized to spawn sub-agents, which in turn spawn further descendants — a single drifted node at depth $d$ with branching factor $n$ can propagate unauthorized scope to $n^d$ descendant agents before the divergence is detected. At $d=10$ and $n=2$, that is 1,024 drifted agents in the field simultaneously. At $n=8$ (a plausible parallelism bound for a well-instrumented agentic runtime), it is over one billion. Drift is not a rare tail event; it is the operative condition of any recursive spawning system that has not been architecturally hardened against it.

The conventional mitigation toolkit for multi-agent authority is thin and inadequate to this problem. Time-to-live (TTL) bounds — the dominant approach in current agentic runtimes — impose a maximum lifetime on each spawned agent, after which it is forcibly terminated. Depth limits restrict the maximum depth of the spawn tree. Neither mechanism addresses the core structural problem: the parent pointer chain that connects an agent to its authorizing human principal can be severed, obscured, or never established in the first place, without any architectural means to detect or prevent it.

The insufficiency of TTL and depth limits is not a matter of degree — it is a matter of kind. A TTL bound prevents an agent from living indefinitely, but it says nothing about whether the agent's current actions are authorized during its lifetime. A depth limit prevents unbounded tree growth, but a drifted agent at depth two with unrestricted scope is no less dangerous than a drifted agent at depth ten. Both mechanisms operate on the \textit{temporal} and \textit{morphological} properties of the spawn tree, not on the \textit{authorization lineage} that determines whether each node is operating within a human-anchored warrant. Drift is a property of the chain, not of its length or duration.

This paper makes the following observation as its core contribution: if the parent pointer is made \textit{immutable} — if the relationship between a child agent and its parent is fixed at spawn time and can never be retroactively modified — then drift becomes structurally impossible. An immutable parent pointer ensures that the chain of delegation is permanently legible. Every agent carries, as a structural property of its existence, a traceable reference to the node that authorized it. That reference cannot be obscured by termination, re-parenting, or scope mutation. The chain is the audit log; the audit log is the chain.

From this core insight, we construct the Recursive Spawning Protocol (RSP), a set of architectural guarantees within the BX3 Framework that enforce immutable parent pointers, mandatory anchor verification at every spawn event, and a Fact Layer pre-commit hook that prevents the creation of agents whose parent cannot be verified against a reachable human principal. RSP operates at the architectural level, not at the agent level — it does not depend on any agent behaving correctly in order for the system to remain accountable. Accountability is a structural fallback, not an emergent property of cooperative agents.

% -----------------------------------------------------------------------
\subsection{Formal Definition of Drift}

We formalize drift as follows. Let $\mathcal{A}$ be the set of active agents. Each agent $a \in \mathcal{A}$ carries a parent pointer $\parent(a)$ referencing the agent that authorized its provisioning; the root of any authorized chain is a human principal $h \in \mathcal{H}$. An agent $a$ is \textit{drifted} at time $t$ when:

\begin{enumerate}
    \item $\status(a, t) = \text{ACTIVE}$ (the agent is executing), and
    \item there exists no path through the parent pointer chain from $a$ to any $h \in \mathcal{H}$ such that each node on that path satisfies $\status(\cdot, t) \neq \text{TERMINATED}$ and $\status(\cdot, t) \neq \text{UNREACHABLE}$.
\end{enumerate}

Drift can arise through several failure modes: the parent terminates silently (orphan drift), the parent redirects its parent pointer to obscure authorization history (re-parenting drift), an agent spawns a child without any verified human anchor (ghost authority drift), or successive scope mutations within an agent cause its effective operating envelope to diverge from what its authorizing human principal sanctioned (scope creep drift). Critically, drift is \textit{heritable}: if agent $a$ is drifted and spawns child $c$, then $c$ is drifted by construction, propagating the condition through the entire descendant subtree.

The drift triad — the conjunction of \textit{orphaned} (parent severed), \textit{drifted} (no human anchor reachable), and \textit{operating} (event loop active) — is the terminal failure condition RSP is designed to make structurally impossible.

% -----------------------------------------------------------------------
\subsection{Paper Roadmap}

The remainder of this paper is organized as follows. Section 2 reviews prior work on multi-agent safety architectures, TTL-based lifecycle management, depth-limited spawning, and the limits of agent-level accountability mechanisms. Section 3 defines drift formally, establishes the failure mode taxonomy, and demonstrates why depth limits alone cannot prevent drift propagation. Section 4 introduces the BX3 Framework's three-layer architecture and shows how RSP is instantiated as an architectural property of the Fact Layer, including the immutable parent pointer invariant, the mandatory anchor verification pre-commit hook, and the scope refinement constraint. Section 5 presents the implementation of RSP within the AgentOS runtime, with live metrics on spawn authorization throughput, anchor verification latency, and zero-drift compliance under adversarial injection conditions. Section 6 evaluates the framework against four drift provocation scenarios and presents comparative results against TTL-baseline systems. Section 7 discusses limitations, including the costs of mandatory anchor verification in high-throughput spawning environments and the open question of anchor re-establishment for long-running agent lineages. Section 8 concludes with a summary of contributions and the research agenda for RSP adoption across Bxthre3 Inc's venture portfolio, including the Irrig8 precision agriculture OS and the Agentic workforce orchestration platform.